\subsection*{为什么平方会带出一个二}
边长为 $x$ 的正方形，每边增加 $h$，新增区域由两条宽为 $h$ 的长条和一个小角组成。两条长条总面积为 $2xh$，小角为 $h^2$。除以边长增量 $h$ 后，长条贡献 $2x$，小角只贡献 $h$，趋零时消失。

这解释了二次幂导数中的系数二来自两条边，不是套公式的巧合。立方体则有三个主要薄层，每层体积约为 $x^2h$；交叠边角含更多个 $h$，在除以 $h$ 后仍趋零，所以得到 $3x^2$。

\subsection*{从图形推广到代数}
更高次幂不容易画，但展开乘积 $(x+h)^n$ 时可以数贡献。一阶项必须从 $n$ 个因子中恰好选一个 $h$，其余选 $x$，一共有 $n$ 种选法，所以系数是 $nx^{n-1}$。选两个及以上 $h$ 的项都带 $h^2$ 或更高次。

“高阶项小”必须说明与什么相比。导数计算比较的是它们相对于 $h$ 的大小，而非只看绝对值。如果一个误差恰为 $3h$，它虽然趋零，除以 $h$ 后仍留下 $3$，不能丢掉。

把 $x=2,h=0.1$ 代入立方增量，先算精确值，再算 $3x^2h$ 的预测，最后减出误差。这个小实验会把导数理解为有误差控制需求的局部预测器。自然数幂的规则可以统一证明，非整数幂则要重新核对定义域。
